\documentclass[preprint,prb,aps,prb]{revtex4}
%\documentclass{revtex4}
%\documentclass[twocolumn,prb,aps,prb]{revtex4}
%\documentstyle[aps,pre,preprint]{revtex}
\usepackage{graphicx}
\usepackage{psfrag}
\usepackage{epsfig}
\usepackage{colordvi}
\usepackage{subfigure}
\usepackage{amsbsy}
\renewcommand{\topfraction}{1}
\renewcommand{\bottomfraction}{1}
\renewcommand{\dbltopfraction}{1}
\renewcommand{\floatpagefraction}{1}
\renewcommand{\textfraction}{0}
\renewcommand{\floatsep}{0pt}
\renewcommand{\dblfloatsep}{0pt}


%\documentclass{article}

%\usepackage[utf8]{inputenc}
%\usepackage[T1]{fontenc}
%\usepackage{lmodern}
%\usepackage{geometry}
%\geometry{verbose,tmargin=2.5cm,bmargin=2.5cm,lmargin=3cm,rmargin=3cm}
%\setcounter{secnumdepth}{3}
%\setcounter{tocdepth}{2}

\usepackage{amsmath}
\usepackage{amssymb}
\usepackage{mathtools}
%\usepackage{xcolor}
%\usepackage{parskip}
%\usepackage[bottom]{footmisc}
%\usepackage{graphicx}
%\usepackage{float}
%\usepackage{subcaption}
\usepackage{ulem}
%\usepackage{hyperref}



%\usepackage[sorting=none, backend=biber]{biblatex}
%\addbibresource{paper_balance.bib}

% This is for authors and affiliation for now
%\usepackage{authblk}

\newcommand{\rd}{{\rm d}}
\newcommand{\bm}{{\mathbf m}}
\newcommand{\bE}{{\mathbf E}}
\newcommand{\bj}{{\mathbf j}}
\newcommand{\br}{{\mathbf r}}
\newcommand{\bB}{{\mathbf B}}
\newcommand{\bh}{{\mathbf h}}
\newcommand{\bk}{{\mathbf k}}
\newcommand{\bV}{{\mathbf V}}
\newcommand{\bv}{{\mathbf v}}
\newcommand{\bR}{{\mathbf R}}
\newcommand{\bA}{{\mathbf A}}
\newcommand{\bJ}{{\mathbf J}}
\newcommand{\bp}{{\mathbf p}}
\newcommand{\ved}{V_{ed}}
\newcommand{\vExB}{V_{E \times B}}
\newcommand{\veperp}{V_{e\perp}}
\newcommand{\pdr}{\partial_r}
\newcommand{\Dqle}{D_{e,22}^{\mathrm{ql}}}
%\newcommand{\Dae}{D_{e22}^{\mathrm{a}}}
\newcommand{\Dae}{D^{\mathrm{a}}}
\newcommand{\Da}{D^{\mathrm{a}}}
\newcommand{\Brplas}{\widetilde B_r^{\mathrm{plas}}}
\newcommand{\Brvac}{\widetilde B_r^{\mathrm{vac}}}
\newcommand{\Irmp}{I_{\mathrm{RMP}}}
\newcommand{\Iexpt}{I_{\mathrm{expt}}}
\newcommand{\qlpartflux}{\Gamma^{\mathrm{(EM)}}}
\newcommand{\qlheatflux}{Q^{\mathrm{(EM)}}}
\newcommand{\Brpert}{\widetilde{B}_r}

\newcommand{\balpha}{{\boldsymbol\alpha}}
\newcommand{\bbeta}{{\boldsymbol\beta}}
\newcommand{\bgamma}{{\boldsymbol\gamma}}
\newcommand{\bdelta}{{\boldsymbol\delta}}
\newcommand{\bepsilon}{{\boldsymbol\epsilon}}
\newcommand{\bvarepsilon}{{\boldsymbol\varepsilon}}
\newcommand{\bzeta}{{\boldsymbol\zeta}}
\newcommand{\boldeta}{{\boldsymbol\eta}}
\newcommand{\btheta}{{\boldsymbol\theta}}
\newcommand{\bvartheta}{{\boldsymbol\vartheta}}
\newcommand{\biota}{{\boldsymbol\iota}}
\newcommand{\bkappa}{{\boldsymbol\kappa}}
\newcommand{\blambda}{{\boldsymbol\lambda}}
\newcommand{\bmu}{{\boldsymbol\mu}}

\newcommand{\bnu}{{\boldsymbol\nu}}
\newcommand{\bxi}{{\boldsymbol\xi}}
\newcommand{\bpi}{{\boldsymbol\pi}}
\newcommand{\bvarpi}{{\boldsymbol\varpi}}
\newcommand{\brho}{{\boldsymbol\rho}}
\newcommand{\bsigma}{{\boldsymbol\sigma}}
\newcommand{\bvarsigma}{{\boldsymbol\varsigma}}
 \newcommand{\btau}{{\boldsymbol\tau}}
 \newcommand{\bupsilon}{{\boldsymbol\upsilon}}
 \newcommand{\bphi}{{\boldsymbol\phi}}
 \newcommand{\bvarphi}{{\boldsymbol\varphi}}
 \newcommand{\bchi}{{\boldsymbol\chi}}
 \newcommand{\bpsi}{{\boldsymbol\psi}}
 \newcommand{\bomega}{{\boldsymbol\omega}}
\newcommand{\bGamma}{{\boldsymbol\Gamma}}
\newcommand{\bDelta}{{\boldsymbol\Delta}}
\newcommand{\bTheta}{{\boldsymbol\Theta}}
\newcommand{\bLambda}{{\boldsymbol\Lambda}}
\newcommand{\bXi}{{\boldsymbol\Xi}}
\newcommand{\bPi}{{\boldsymbol\Pi}}
\newcommand{\bSigma}{{\boldsymbol\Sigma}}
\newcommand{\bUpsilon}{{\boldsymbol\Upsilon}}
\newcommand{\bPhi}{{\boldsymbol\Phi}}
\newcommand{\bPsi}{{\boldsymbol\Psi}}
\newcommand{\bOmega}{{\boldsymbol\Omega}}

\newcommand{\cA}{{\mathcal A}}
\newcommand{\cB}{{\mathcal B}}
\newcommand{\cC}{{\mathcal C}}
\newcommand{\cD}{{\mathcal D}}
\newcommand{\cE}{{\mathcal E}}
\newcommand{\cF}{{\mathcal F}}
\newcommand{\cG}{{\mathcal G}}
\newcommand{\cH}{{\mathcal H}}
\newcommand{\cI}{{\mathcal I}}
\newcommand{\cJ}{{\mathcal{J}}}
\newcommand{\cK}{{\mathcal K}}
\newcommand{\cL}{{\mathcal L}}
\newcommand{\cM}{{\mathcal M}}
\newcommand{\cN}{{\mathcal N}}
\newcommand{\cO}{{\mathcal O}}
\newcommand{\cP}{{\mathcal P}}
\newcommand{\cQ}{{\mathcal Q}}
\newcommand{\cR}{{\mathcal R}}
\newcommand{\cS}{{\mathcal S}}
\newcommand{\cT}{{\mathcal T}}
\newcommand{\cU}{{\mathcal U}}
\newcommand{\cV}{{\mathcal{V}}}
\newcommand{\cW}{{\mathcal W}}
\newcommand{\cX}{{\mathcal X}}
\newcommand{\cY}{{\mathcal Y}}
\newcommand{\cZ}{{\mathcal Z}}


\newcommand{\be}[1]{\begin{equation} \label{#1}}
\newcommand{\ee}{\end{equation}}
\newcommand{\bea}[1]{\begin{eqnarray} \label{#1}}
\newcommand{\eea}{\end{eqnarray}}
\newcommand{\bean}{\begin{eqnarray*}}
\newcommand{\eean}{\end{eqnarray*}}

\newcommand{\non}{\nonumber\\}
\newcommand{\eq}[1]{(\ref{#1})}
\newcommand{\difp}[2]{\frac{\partial #1}{\partial #2}}

% For editing
\usepackage{hyphenat}
\usepackage{ifthen}
\usepackage{comment}

%\ifthenelse{1=1} %change to 1=2 for final version without comments
%{
%    \specialcomment{todo}{\begingroup\ttfamily\footnotesize\color{red}TODO: }{\par\endgroup}
%    \specialcomment{remark}{\begingroup\ttfamily\footnotesize\color{blue}REMARK: }{\endgroup}
%}
%{
%    \excludecomment{todo}
%    \excludecomment{remark}
%}

\begin{document}

\section{Transport due to electrostatic perturbations}

For the treatment of electron component, its is sufficient to use the drift-kinetic approximation.
Drit kinetic equation using for velocity space variables besides the gyrophase $\phi$ also
the magnetic moment $\mu$ and parallel velocity $v_\parallel$ where
$$
\mu = \frac{m_\alpha v_\perp^2}{2 B}=\frac{e_\alpha  J_\perp}{m_\alpha  c}
$$
and $J_\perp$ is the perpendicular adiabatic invariant,
is of the form
\be{drikineq}
\difp{f}{t}+\bv_g\cdot\nabla f + \dot v_\parallel \difp{f}{v_\parallel}= \hat L_C f.
\ee
Here, the guiding center velocity is
\be{gvcvelpp}
\bv_g = \frac{v_\parallel \bB}{B_\parallel^\ast} 
+ \frac{\bB\times\nabla \left(\mu B + e_\alpha \Phi\right)}
{m_\alpha \omega_c B_\parallel^\ast}
+\frac{v_\parallel^2}{B_\parallel^\ast}\nabla\times \frac{\bB}{\omega_c},
\qquad
B_\parallel^\ast = B + v_\parallel \bh\cdot \nabla\times\frac{\bB}{\omega_c},
\ee
$\bh = \bB/B$ and $\omega_c=e_\alpha  B/(m_\alpha c)$ are unit vector along the magnetic field and the cyclotron frequency, respectively,
\be{vpardot}
\dot v_\parallel = - \frac{1}{m_\alpha  B_\parallel^\ast}\left(\bB+ v_\parallel\nabla\times \frac{\bB}{\omega_c}\right)
\cdot\nabla \left(\mu B + e_\alpha \Phi\right),
\ee
and $\hat L_C$ is the collision opertor.
We ignore the magnetic drift in favour of electric drift so that $B_\parallel^\ast \rightarrow B$ and
\be{nomagdrift}
\bv_g \rightarrow v_\parallel \bh + \frac{c \bh\times\nabla\Phi}{B},
\qquad
\dot v_\parallel \rightarrow - \frac{1}{m_\alpha } \bh\cdot\nabla\left(\mu B + e_\alpha \Phi\right). 
\ee

\subsection{Quasilinear theory}

We assume purely electrostatic perturbations such that magnetic field $\bB$ 
is axisymmetric and electrostatic potential $\Phi=\Phi_0 + \delta \Phi$ contains
the unperturbed, axisymmetric potential $\Phi_0=\Phi_0(r)$ and the perturbation 
$\delta \Phi$.
We stay in a straight cylinder geometry where the unperturbed motion is simple,
\be{phsp_vel_0}
\bv_{g0} = v_\parallel \bh+\bV_E,
\qquad
\bV_E = \frac{c \Phi_0^\prime(r)}{B}\bh\times\nabla r,
\qquad
\dot v_{\parallel 0} = 0.
\ee
Respectively, the perturbation of the phase space velocity $\delta \bv_{g} = \bv_{g}-\bv_{g0}$ is driven by $\delta \Phi$,
\be{phsp_vel_1}
\delta \bv_{g} = \frac{c \bh\times\nabla\delta\Phi}{B},
\qquad
\delta\dot v_{\parallel} = - \frac{e_\alpha}{m_\alpha } \bh\cdot\nabla\Phi.
\ee

We use the normalized coordinates $(r,\vartheta,\varphi)$ where $\varphi=z/R_0$ is the normalized cylinder length
and $R_0 = const$ is the major radius, so that $\sqrt{g} = r R_0$.
In zero order over the perturbation  amplitude, we take a local Maxwellian,
\be{boltz}
f_0(r,\mu,v_\parallel)=
\frac{n_\alpha(r)}{(2\pi m_\alpha  T_\alpha(r))^{3/2}}
\exp\left(-\frac{2\mu B(r) + m_\alpha v_\parallel^2}{2 T_\alpha(r)}\right),
\ee
which obviously satisfies the unperturbed steady state equation
$$
\bv_{g0} \cdot \nabla f_0 + \dot v_{\parallel 0} \difp{f_0}{v_\parallel} = \hat L_C f_0.
$$
We assume the perturbation of the form
\be{pertPhi_form}
\delta \Phi(t,r,\vartheta,\varphi) = {\rm Re}\sum_\bm \Phi_\bm(r) {\rm e}^{im\vartheta+i n \varphi -i \omega_\bm t},
\ee
where $\bm=(m,n)$. We similarly expand the perturbation of the distribution function, $\delta f = f-f_0$, which we treat in the linear order (omitting the terms quadratic in $\delta\Phi$),
\be{pertf_form}
\delta f(t,r,\vartheta,\varphi,\mu,v_\parallel) = {\rm Re} \sum_\bm f_\bm(r,\mu,v_\parallel) 
{\rm e}^{im\vartheta+i n \varphi -i \omega_\bm t},
\ee
and obtain equations for Fourier amplitudes,
\be{eqforamps}
\left(\bk\cdot \bv_{g0}- \omega_\bm\right) f_\bm + i \hat L_C f_\bm 
= 
\left(
\left(\frac{n_\alpha^\prime}{n_\alpha}-\frac{3 T_\alpha^\prime}{2 T_\alpha}
+\frac{T_\alpha^\prime}{T_\alpha}\frac{m_\alpha v^2}{2 T_\alpha}\right)
\frac{ck_\perp}{B}
-\frac{e_\alpha k_\parallel v_\parallel}{T_\alpha}
\right)
f_0  \Phi_\bm,
\ee
where we denoted
\be{kvecs}
\bk = m\nabla\vartheta+n\nabla\varphi,
\qquad
k_\perp = \bk\cdot\bh\times\nabla r,
\qquad
k_\parallel = \bk\cdot\bh,
\ee
and ignored in the r.h.s. term with $B^\prime$ which enters $n_\alpha^\prime / n_\alpha - \mu B^\prime/T_\alpha + \dots$
for the same reason as the magnetic drift. We have also denoted the kinetic energy as
$$
\frac{m_\alpha v^2}{2} = \mu B +\frac{m_\alpha v_\parallel^2}{2}.
$$
We consider high frequency short scale perturbations such that 
${\rm max}(\omega_\bm, k_\perp V_E) \gg \nu$ where $\nu$ is the collision frequency.
For these perturbations, transport is determined by the ``collisionless'' (Landau damping) mechanism such that it
is sufficient to use Krook collision model, $\hat L_C = - \nu$, because small $\nu$ does not enter the result.
Thus, Eq.~\eq{eqforamps} becomes algebraic,
\be{eqforamps_alg}
\left(k_\parallel v_\parallel + k_\perp V_E-\omega_\bm-i\nu\right) f_\bm 
=
\left(
\left(\frac{n_\alpha^\prime}{n_\alpha}-\frac{3 T_\alpha^\prime}{2 T_\alpha}
+\frac{T_\alpha^\prime}{T_\alpha}\frac{m_\alpha v^2}{2 T_\alpha}\right)
\frac{ck_\perp}{B}
-\frac{e_\alpha k_\parallel v_\parallel}{T_\alpha}
\right)
f_0  \Phi_\bm,
\ee
where
\be{VE}
V_E = \frac{c \Phi_0^\prime}{B}.
\ee
Multiplying the perturbation of the distribution function~\eq{pertf_form}
with radial guding center velocity $\delta \bv_g\cdot\nabla r$ where $\delta \bv_g$ is
given by Eq.~\eq{phsp_vel_1}, averaging the result over the angles and integrating over the momentum 
space, we obtain particle flux density as
\bea{Gamma}
\Gamma_\alpha
&=& 
\int\rd^3 p f_0 \frac{1}{2}
{\rm Re}
\sum_\bm 
\frac{i c^2 k_\perp^2 \left|\Phi_\bm\right|^2}{B^2\left(k_\parallel v_\parallel + k_\perp V_E-\omega_\bm-i\nu\right)}
\left(
\frac{n_\alpha^\prime}{n_\alpha}-\frac{3 T_\alpha^\prime}{2 T_\alpha}
+\frac{T_\alpha^\prime}{T_\alpha}\frac{m_\alpha v^2}{2 T_\alpha}
-\frac{e_\alpha k_\parallel v_\parallel B}{T_\alpha ck_\perp}
\right)
\nonumber \\
&\rightarrow &
-\frac{\pi}{2}\int\rd^3 p f_0 
\sum_\bm
\frac{c^2 k_\perp^2 \left|\Phi_\bm\right|^2}{B^2}
\delta\left(k_\parallel v_\parallel + k_\perp V_E-\omega_\bm\right)
\left(
\frac{n_\alpha^\prime}{n_\alpha}-\frac{3 T_\alpha^\prime}{2 T_\alpha}
+\frac{T_\alpha^\prime}{T_\alpha}\frac{m_\alpha v^2}{2 T_\alpha}
-\frac{e_\alpha k_\parallel v_\parallel B}{T_\alpha ck_\perp}
\right)
\nonumber \\
&=&
-\frac{\pi}{2}\int\rd^3 p f_0
\sum_\bm
\frac{c^2 k_\perp^2 \left|\Phi_\bm\right|^2}{B^2}
\delta\left(k_\parallel v_\parallel + k_\perp V_E-\omega_\bm\right)
\left(A_1
+\frac{m_\alpha v^2}{2 T_\alpha} A_2
-\frac{\omega_\bm}{k_\perp D_B}
\right),
\eea
where we introduced the following notation for the thermodynamic forces and Bohm diffusion coefficient,
\be{A_1_A_2_D_B}
A_1 = \frac{n_\alpha^\prime}{n_\alpha}+\frac{e_\alpha \Phi_0^\prime}{T_\alpha}-\frac{3 T_\alpha^\prime}{2 T_\alpha},
\qquad
A_2 = \frac{T_\alpha^\prime}{T_\alpha},
\qquad
D_B = \frac{c T_\alpha}{e_\alpha B}.
\ee
Integration over the momentum space,
$$
\int \rd^3 p = 2\pi m_\alpha^2 B\int\limits_0^\infty \rd \mu\int\limits_{-\infty}^\infty \rd v_\parallel,
$$
is trivial over $\mu$ and easy over $v_\parallel$ where the $\delta$-function should be used to give
\bea{Gamma_triv}
\Gamma_\alpha
&=&
- n_\alpha \left(\frac{\pi m_\alpha}{8 T_\alpha}\right)^{1/2}
\int\limits_{-\infty}^\infty \rd v_\parallel
\exp\left(- \frac{m_\alpha v_\parallel^2}{2 T_\alpha}\right)
\sum_\bm
\frac{c^2 k_\perp^2 \left|\Phi_\bm\right|^2}{B^2}
\delta\left(k_\parallel v_\parallel + k_\perp V_E-\omega_\bm\right)
\nonumber \\
&\times&
\left(A_1+A_2
+\frac{m_\alpha v_\parallel^2}{2 T_\alpha} A_2
-\frac{\omega_\bm}{k_\perp D_B}
\right)
\nonumber \\
&=&
- n_\alpha \left(\frac{\pi m_\alpha}{8 T_\alpha}\right)^{1/2}
\sum_\bm
\frac{c^2 k_\perp^2 \left|\Phi_\bm\right|^2}{|k_\parallel| B^2}
\exp\left(- \frac{m_\alpha \left(\omega_\bm-k_\perp V_E\right)^2}{2 k_\parallel^2 T_\alpha}\right)
\nonumber \\
&\times&
\left(A_1+A_2
+
\frac{m_\alpha \left(\omega_\bm-k_\perp V_E\right)^2}{2 k_\parallel^2 T_\alpha}
A_2
-\frac{\omega_\bm}{k_\perp D_B}
\right)
\nonumber \\
&\equiv&
- n_\alpha (D_{11}^\alpha A_1 + D_{12}^\alpha A_2 - V^\alpha_1),
\eea
where the last expression defines diffusion coefficients and convection velocity.
In particular, particle diffusion coeficient follows from this expression as
\be{D11}
D_{11}^\alpha = \left(\frac{\pi m_\alpha}{8 T_\alpha}\right)^{1/2}
\sum_\bm
\frac{c^2 k_\perp^2 \left|\Phi_\bm\right|^2}{|k_\parallel| B^2}
\exp\left(- \frac{m_\alpha \left(\omega_\bm-k_\perp V_E\right)^2}{2 k_\parallel^2 T_\alpha}\right).
\ee

\subsection{Static noise on the GORILLA grid}

For numerical experiments in GORILLA, we generate the potential perturbation by adding in each grid node $k$
a random potential value using the standard random numbers $\xi$,
\be{delphik}
\delta \Phi_k = \Phi_A \xi_k, 
\qquad
\overline{\xi_k}=0,
\qquad
\overline{\xi_k \xi_{k^\prime}}=\delta_{k,k^\prime}.
\ee
In a 1D exapmple, potential perturbation is then given by the continuous piecewise-linear periodic function 
of an angle $\delta \Phi(\vartheta)$. Let us compute its Fourier amplitude,
\be{fourampnoise}
\Phi_m 
= 
\frac{1}{2\pi}\int\limits_0^{2\pi}\rd \vartheta {\rm e}^{-im\vartheta}\delta \Phi(\vartheta)
= 
-\frac{1}{2\pi m^2}\int\limits_0^{2\pi}\rd \vartheta {\rm e}^{-im\vartheta}\delta \Phi^{\prime\prime}(\vartheta).
\ee
Second derivative $\delta \Phi^{\prime\prime}(\vartheta)$ 
of the continuous piecewise-linear function $\delta \Phi(\vartheta)$ 
equals to the sum of $\delta$-functions at the grid
nodes $\vartheta = 2\pi k/K$ where $K$ is grid size and $1\le k \le K$,
\be{deltafuns}
\delta \Phi^{\prime\prime}(\vartheta) = \frac{K\Phi_A}{2\pi}\sum\limits_{k=1}^K
\left(\xi_{k+1}+\xi_{-1}-2\xi_{k}\right)
\delta\left(\vartheta - \frac{2\pi k}{K}\right).
\ee
Respectively, Fourier amplitude~\eq{fourampnoise} is
\be{deltafuns_four}
\Phi_m
= 
-\frac{K\Phi_A}{(2\pi m)^2}\sum\limits_{k=1}^K
\left(\xi_{k+1}+\xi_{k-1}-2\xi_{k}\right)
\ee
Obviously, its expectation value is zero, $\overline{\Phi_m}=0$. We are interested, however, in the expectation values
of its module square,
\be{modsqexpect}
\overline{|\Phi_m|^2} = \frac{K^2\Phi_A^2}{(2\pi m)^4} \sum\limits_{k,k^\prime=1}^K
\overline{\left(\xi_{k+1}+\xi_{k-1}-2\xi_{k}\right)\left(\xi_{k^\prime+1}+\xi_{k^\prime-1}-2\xi_{k^\prime}\right)}
\exp\left(\frac{2\pi i m (k-k^\prime)}{K}\right).
\ee
Using the properties~\eq{delphik} of random numbers we get
$$
\overline{\left(\xi_{k+1}+\xi_{k-1}-2\xi_{k}\right)\left(\xi_{k^\prime+1}+\xi_{k^\prime-1}-2\xi_{k^\prime}\right)}
=
6 \delta_{k,k^\prime}-4 \delta_{k-1,k^\prime}-4 \delta_{k+1,k^\prime}+\delta_{k-2,k^\prime}+\delta_{k+2,k^\prime},
$$
so that~\eq{modsqexpect} becomes
\be{modsqexpect_fin}
\overline{|\Phi_m|^2} = \frac{K^3\Phi_A^2}{(2\pi m)^4} 
\left(6-8\cos\left(\frac{2\pi m}{K}\right)+2\cos\left(\frac{4\pi m}{K}\right)\right).
\ee
I am lazy to compute this for a 2D grid in $(\vartheta,\varphi)$ variables with the size $K\times L$. Therefore
I write the respectve fomula without a proof,
\bea{modsqexpect_fin2D}
\overline{|\Phi_\bm|^2} &=& \frac{K^3L^3\Phi_A^2}{(2\pi m)^4 (2\pi n)^4} 
\left(6-8\cos\left(\frac{2\pi m}{K}\right)+2\cos\left(\frac{4\pi m}{K}\right)\right)
\nonumber \\
&\times& \qquad\qquad\qquad \left(6-8\cos\left(\frac{2\pi n}{L}\right)+2\cos\left(\frac{4\pi n}{L}\right)\right).
\eea
Let us estimate a diffusion coefficient~\eq{D11} using expectation value~\eq{modsqexpect_fin2D} 
for the amplitude $|\Phi_\bm|^2$ setting $\omega_\bm=0$ becuse of static perturbations and
setting also $V_E=0$. The latter setting corresponds to zero equilibrium electric field, however even
for the finite but reasonable electric field this term is of the order of 
$V_E^2 k_\perp^2/(v_{T\alpha}^2 k_\parallel^2) \sim (k_\perp^2 \rho_{L\alpha^2})/(k_\parallel^2 r^2) \ll 1$
for $k_\parallel \sim k_\perp$. Thus we set the exponent in~\eq{D11} to one. Approximating roughly
$k_\parallel \approx n/R_0$ and $k_\perp \approx m/r$ we get
\bea{D11_grid}
D_{11}^\alpha 
&\approx& 
\left(\frac{\pi m_\alpha}{8 T_\alpha}\right)^{1/2}
\frac{c^2 R_0 \Phi_A^2}{r^2 B^2}\frac{1}{KL}
\sum_\bm
\frac{m^2}{|n|}\left(\frac{2\pi m}{K}\right)^{-4} \left(\frac{2\pi n}{L}\right)^{-4}
\\
&\times&
\left(6-8\cos\left(\frac{2\pi m}{K}\right)+2\cos\left(\frac{4\pi m}{K}\right)\right)
\left(6-8\cos\left(\frac{2\pi n}{L}\right)+2\cos\left(\frac{4\pi n}{L}\right)\right).
\nonumber
\eea
since $2\pi/K \ll 1$ and $2\pi/L \ll 1$, we replace summation by the integration introducing
the variables $x = 2\pi m/K$ and $y=2\pi n/L$ so that
$$
\frac{1}{KL}\sum_\bm \approx \frac{1}{4\pi^2}
\int\limits_{-\infty}^\infty \rd x
\int\limits_{-\infty}^\infty \rd y,
$$
and
\bea{D11_grid_cont}
D_{11}^\alpha
&\approx&
\left(\frac{\pi m_\alpha}{8 T_\alpha}\right)^{1/2}
\frac{c^2 R_0 \Phi_A^2}{r^2 B^2}
\frac{K^2}{8\pi^3 L}
\int\limits_{-\infty}^\infty \rd x
\int\limits_{-\infty}^\infty \rd y
\frac{1}{x^2 |y|^5}
\\
&\times&
\left(6-8\cos\left(x\right)+2\cos\left(2x\right)\right)
\left(6-8\cos\left(y\right)+2\cos\left(2y\right)\right).
\nonumber
\eea
We note that second line of Eq.~\eq{D11_grid_cont} scales at small $x$ and $y$ as $x^4 y^4$.
Thus, integral over $x$ strictly converges giving
$$
\int\limits_{-\infty}^\infty \rd x \frac{6-8\cos\left(x\right)+2\cos\left(2x\right)}{x^2}  
=
\int\limits_{-\infty}^\infty \rd x \frac{8\sin\left(x\right)-4\sin\left(2x\right)}{x}
=
4 \int\limits_{-\infty}^\infty \rd x \frac{\sin x}{x}
= 4 \pi
$$
but integral over $y$ logarithmically diverges.
This divergence however is unphysical because it is the result of using collisionless limit for
low $k_\parallel$ within the assumption $k_\parallel v_\parallel \gg \nu$. We can introduce some minimum
value of $y$
$$
\int\limits_{-\infty}^\infty \rd y \frac{6-8\cos\left(y\right)+2\cos\left(2y\right)}{|y|^5}
\rightarrow
\int\limits_{y_\text{min}}^\infty \rd y \frac{6-8\cos\left(y\right)+2\cos\left(2y\right)}{y^5}
\sim
\int\limits_{y_\text{min}}^1 \rd y \frac{1}{y} = \left|\log(y_\text{min})\right|.
$$
Thus we get an estimate
\be{D11_grid_cont_est}
D_{11}^\alpha
\sim
\frac{\left|\log(y_\text{min})\right|}{2^{5/2}\pi^{3/2}}
\left(\frac{m_\alpha}{T_\alpha}\right)^{1/2}
\frac{c^2 R_0 \Phi_A^2}{r^2 B^2}
\frac{K^2}{L}.
\ee


\end{document}
